\paragraph{环边积分分解与 Vieta 系数闭式：例外幂和的可枚举结构化重写}
沿用本节记号：\(q\ge 2\)、\(U_q=\{0,1\}^q\)、\(T^{(q)}\) 为软核矩阵，
\(\{\lambda_i^{(q)}\}_{i=0}^{q}\) 为 \(S_q\)-对称子空间上的 \(q+1\) 个单根例外特征值，
\[
S_m(q):=\sum_{i=0}^{q}\bigl(\lambda_i^{(q)}\bigr)^m,\qquad m\ge 1.
\]

\begin{conclusion}[例外幂和的环边路径积分分解与 Fibonacci 组件因子化]\label{con:xi-terminal-replica-softcore-cycle-edge-integral-decomposition-fibonacci-components}
固定 \(m\ge 1\)。记环的边位集合
\[
E_m:=\{(t,t+1):t\in\ZZ/m\ZZ\},
\]
并对任意 \(S\subseteq E_m\) 定义
\[
N_m(S):=\#\left\{x\in\{0,1\}^{\ZZ/m\ZZ}:\
x_t x_{t+1}=0\ \text{对全部 }(t,t+1)\in S\right\}.
\]
再定义
\[
\Omega_m(q):=\sum_{j=0}^{q}(-1)^{jm}\varphi^{m(q-2j)}.
\]
则对一切 \(q\ge 2\)、\(m\ge 1\)，有
\begin{equation}\label{eq:xi-cycle-edge-integral-trace}
\Tr\!\bigl((T^{(q)})^m\bigr)=2^{-m}\sum_{S\subseteq E_m}N_m(S)^q,
\end{equation}
\begin{equation}\label{eq:xi-cycle-edge-integral-exceptional}
S_m(q)=2^{-m}\!\left(\Omega_m(q)+\sum_{S\subsetneq E_m}N_m(S)^q\right).
\end{equation}
此外，当 \(S\subsetneq E_m\) 时，关联无向图分解为若干路径分量；若其分量顶点数为
\(
n_1,\dots,n_r
\)
（\(\sum_{\ell=1}^{r}n_\ell=m\)），则
\begin{equation}\label{eq:xi-cycle-edge-integral-fibonacci-factor}
N_m(S)=\prod_{\ell=1}^{r}F_{n_\ell+2}.
\end{equation}
\end{conclusion}

\begin{proof}
\textbf{第一步：迹矩展开与坐标解耦。}
由
\(
T^{(q)}_{u,v}=\frac12(1+\mathbf 1_{\{u\wedge v=0\}})
\)
及环路径迹展开，
\[
\Tr\!\bigl((T^{(q)})^m\bigr)
=\sum_{u_0,\dots,u_{m-1}\in U_q}
\prod_{t\in\ZZ/m\ZZ}T^{(q)}_{u_t,u_{t+1}},\qquad u_m:=u_0.
\]
对 \(m\) 个因子完全展开得
\[
\Tr\!\bigl((T^{(q)})^m\bigr)
=2^{-m}\sum_{S\subseteq E_m}
\#\left\{(u_t)_t\in U_q^m:\ u_t\wedge u_{t+1}=0,\ \forall (t,t+1)\in S\right\}.
\]
固定 \(S\) 后，约束按坐标 \(r=1,\dots,q\) 完全解耦；每个坐标都是同一二元环约束计数 \(N_m(S)\)。
故内层计数等于 \(N_m(S)^q\)，即得 \eqref{eq:xi-cycle-edge-integral-trace}。

\textbf{第二步：例外幂和替换。}
由全谱分解（例外谱与固定谱并置）：
\[
\Tr\!\bigl((T^{(q)})^m\bigr)
=S_m(q)+\sum_{j=0}^{q}\left(\binom{q}{j}-1\right)
\left(\frac{(-1)^j\varphi^{q-2j}}{2}\right)^m.
\]
化简第二项：
\[
2^{-m}\!\left(
\sum_{j=0}^{q}\binom{q}{j}(-1)^{jm}\varphi^{m(q-2j)}
-\sum_{j=0}^{q}(-1)^{jm}\varphi^{m(q-2j)}
\right)
=2^{-m}\!\left(L_m^q-\Omega_m(q)\right).
\]
其中第一和式由二项式定理给出
\[
\sum_{j=0}^{q}\binom{q}{j}(\varphi^m)^{q-j}\bigl((-1)^m\varphi^{-m}\bigr)^j
=\bigl(\varphi^m+(-1)^m\varphi^{-m}\bigr)^q=L_m^q.
\]
故
\[
\Tr\!\bigl((T^{(q)})^m\bigr)=S_m(q)+2^{-m}\bigl(L_m^q-\Omega_m(q)\bigr).
\]
再与 \eqref{eq:xi-cycle-edge-integral-trace} 联立，并用 \(N_m(E_m)=L_m\)（环 \(C_m\) 的独立集数），即得
\eqref{eq:xi-cycle-edge-integral-exceptional}。

\textbf{第三步：真子图的 Fibonacci 乘法分解。}
当 \(S\subsetneq E_m\) 时，去掉至少一个环边位，关联图无环且各分量为路径。
对 \(n\) 顶点路径 \(P_n\)，独立集数 \(a_n\) 满足
\(
a_n=a_{n-1}+a_{n-2}
\)、\(a_0=1,a_1=2\)，故 \(a_n=F_{n+2}\)。
各分量独立相乘得到 \eqref{eq:xi-cycle-edge-integral-fibonacci-factor}。证毕。
\end{proof}

\begin{corollary}[二次幂和的完全闭式与 \texorpdfstring{$F_{2q+2}$}{F\_(2q+2)} 出现]\label{cor:xi-terminal-replica-softcore-exceptional-power-sum-m2-fibonacci-closed}
\leanverified{paper\_xi\_terminal\_replica\_softcore\_exceptional\_power\_sum\_m2\_fibonacci\_closed}
对任意 \(q\ge 2\)，有
\begin{equation}\label{eq:xi-exceptional-s2-complete-closed-form}
S_2(q)=\frac14\left(4^q+2\cdot 3^q+F_{2q+2}\right).
\end{equation}
\end{corollary}

\begin{proof}
在结论 \ref{con:xi-terminal-replica-softcore-cycle-edge-integral-decomposition-fibonacci-components} 中取 \(m=2\)。
此时 \(E_2\) 含两个边位；真子集为 \(\varnothing\) 与两个单边位子集，故
\[
N_2(\varnothing)=4,\qquad N_2(S)=3\ \ (|S|=1).
\]
又
\[
\Omega_2(q)=\sum_{j=0}^{q}\varphi^{2(q-2j)}
=\frac{\varphi^{2(q+1)}-\varphi^{-2(q+1)}}{\sqrt5}
=F_{2q+2}.
\]
代入 \eqref{eq:xi-cycle-edge-integral-exceptional} 即得 \eqref{eq:xi-exceptional-s2-complete-closed-form}。证毕。
\end{proof}

\begin{conclusion}[例外特征多项式二阶 Vieta 系数闭式]\label{con:xi-terminal-replica-softcore-exceptional-vieta-e2-closed-form}
记
\[
\chi_{\mathrm{exc}}^{(q)}(x)
:=\prod_{i=0}^{q}\bigl(x-\lambda_i^{(q)}\bigr)
=x^{q+1}-e_1^{(q)}x^q+e_2^{(q)}x^{q-1}-\cdots+(-1)^{q+1}e_{q+1}^{(q)}.
\]
则对任意 \(q\ge 2\)，
\begin{equation}\label{eq:xi-exceptional-vieta-e2-closed}
e_2^{(q)}
=2^{q-2}F_{q+1}
-\frac14\Bigl(3^q+F_qF_{q+1}\Bigr).
\end{equation}
\end{conclusion}

\begin{proof}
由 Vieta 关系，
\[
e_1^{(q)}=\sum_{i=0}^{q}\lambda_i^{(q)}=S_1(q),\qquad
e_2^{(q)}=\sum_{0\le i<j\le q}\lambda_i^{(q)}\lambda_j^{(q)}
=\frac{S_1(q)^2-S_2(q)}{2}.
\]
由结论
\ref{con:xi-terminal-replica-softcore-exceptional-power-sum-m1-integer-collapse}
与推论 \ref{cor:xi-terminal-replica-softcore-exceptional-power-sum-m2-fibonacci-closed}，
\[
S_1(q)=2^{q-1}+\frac12F_{q+1},\qquad
S_2(q)=\frac14\bigl(4^q+2\cdot 3^q+F_{2q+2}\bigr).
\]
代入并用恒等式
\(
F_{2q+2}=F_{q+1}^2+2F_qF_{q+1}
\)
化简，即得 \eqref{eq:xi-exceptional-vieta-e2-closed}。证毕。
\end{proof}

\begin{conclusion}[六次幂和全闭式与 Fibonacci 组件谱系]\label{con:xi-terminal-replica-softcore-exceptional-power-sum-m6-fibonacci-component-spectrum}
对任意 \(q\ge 2\)，
\begin{equation}\label{eq:xi-exceptional-s6-fibonacci-spectrum-closed}
\begin{aligned}
S_6(q)=\frac1{64}\Big(
&64^q+6\cdot48^q+6\cdot40^q+9\cdot36^q+6\cdot32^q+12\cdot30^q\\
&+2\cdot27^q+6\cdot26^q+3\cdot25^q+6\cdot24^q+6\cdot21^q+\Omega_6(q)
\Big),
\end{aligned}
\end{equation}
且 \(\Omega_6\) 满足
\begin{equation}\label{eq:xi-omega6-second-order-recurrence}
\Omega_6(0)=1,\quad
\Omega_6(1)=18,\quad
\Omega_6(q)=18\,\Omega_6(q-1)-\Omega_6(q-2)\ \ (q\ge 2).
\end{equation}
\end{conclusion}

\begin{proof}
由 \eqref{eq:xi-cycle-edge-integral-exceptional}（\(m=6\)）：
\[
S_6(q)=2^{-6}\!\left(\Omega_6(q)+\sum_{S\subsetneq E_6}N_6(S)^q\right).
\]
按缺边位数 \(r:=|E_6\setminus S|\) 分类，结合
\eqref{eq:xi-cycle-edge-integral-fibonacci-factor} 逐类得到 \(N_6(S)\) 与重数：
\begin{enumerate}
  \item \(r=6\)：\(N=64\)，重数 \(1\)；
  \item \(r=5\)：\(N=48\)，重数 \(6\)；
  \item \(r=4\)：\(N=40\)（重数 \(6\)），\(N=36\)（重数 \(9\)）；
  \item \(r=3\)：\(N=32\)（重数 \(6\)），\(N=30\)（重数 \(12\)），\(N=27\)（重数 \(2\)）；
  \item \(r=2\)：\(N=26\)（重数 \(6\)），\(N=24\)（重数 \(6\)），\(N=25\)（重数 \(3\)）；
  \item \(r=1\)：\(N=21\)，重数 \(6\)。
\end{enumerate}
总重数 \(1+6+15+20+15+6=63=2^6-1\)。
代回即得 \eqref{eq:xi-exceptional-s6-fibonacci-spectrum-closed}。

又
\[
\Omega_6(q)=\sum_{j=0}^{q}\varphi^{6(q-2j)}
\]
是特征根 \(\varphi^6,\varphi^{-6}\) 的二指数线性组合，故满足
\[
\Omega_6(q)-(\varphi^6+\varphi^{-6})\Omega_6(q-1)+\Omega_6(q-2)=0.
\]
而 \(\varphi^6+\varphi^{-6}=L_6=18\)，并由定义得初值 \(1,18\)，即 \eqref{eq:xi-omega6-second-order-recurrence}。证毕。
\end{proof}

\begin{proposition}[固定 \texorpdfstring{$m$}{m} 的 \texorpdfstring{$q$}{q}-方向有限阶线性递推]\label{prop:xi-terminal-replica-softcore-fixed-m-q-recurrence-PAQ}
固定 \(m\ge 1\)。定义
\[
\mathcal A_m:=\{N_m(S):S\subsetneq E_m\}\subset\ZZ_{>0},
\]
\[
P_m(t):=\prod_{a\in\mathcal A_m}(t-a),
\qquad
Q_m(t):=t^2-L_m t+(-1)^m.
\]
记移位算子 \(E\) 作用为 \((Ef)(q)=f(q+1)\)。则
\begin{equation}\label{eq:xi-fixed-m-q-recurrence-PAQ}
P_m(E)\,Q_m(E)\,S_m(\cdot)\equiv 0.
\end{equation}
换言之，\(q\mapsto S_m(q)\) 满足一个由 \(\mathcal A_m\) 与 \(L_m\) 显式决定的整系数齐次递推；
其阶数为 \(|\mathcal A_m|+2\)。
\end{proposition}

\begin{proof}
由 \eqref{eq:xi-cycle-edge-integral-exceptional}，
\[
2^mS_m(q)=\Omega_m(q)+\sum_{S\subsetneq E_m}N_m(S)^q.
\]
将真子集项按同值 \(a\in\mathcal A_m\) 归并，写为
\[
\sum_{S\subsetneq E_m}N_m(S)^q=\sum_{a\in\mathcal A_m}c_m(a)\,a^q,\qquad c_m(a)\in\NN.
\]
每个指数序列 \(a^q\) 都被 \(E-a\) 湮灭，故总和被 \(P_m(E)\) 湮灭。

另一方面，\(\Omega_m(q)\) 的特征根为 \(\varphi^m\) 与 \((-1)^m\varphi^{-m}\)，因此满足
\[
\Omega_m(q+2)-L_m\Omega_m(q+1)+(-1)^m\Omega_m(q)=0,
\]
即 \(Q_m(E)\Omega_m\equiv 0\)。
对
\(
2^mS_m=\Omega_m+\sum_{a}c_m(a)a^q
\)
左作用 \(P_m(E)Q_m(E)\) 即得
\eqref{eq:xi-fixed-m-q-recurrence-PAQ}。证毕。
\end{proof}

\endinput
